% ED26-SRC-063-2001/2002: complete Marenbon Fordham manuscript chapters.
% Final2008 published0201/0202 and Wilks's Abelard chapter remain distinct.
\chapter{Early Latin Logic: Arguments, Categories, and the Turn of the Twelfth Century}
\label{sec:p13-early-medieval-latin-logic}
\chaptermark{Early Latin Logic}
\index{medieval logic!early Latin}\index{Marenbon, John}
\index{Boethius}\index{Anselm}\index{topical argument}

The early Latin tradition transmitted more than a curtailed Greek syllabus.
It joined categorical and hypothetical argument, topical discovery, theories
of words and thoughts, and the use of Categories in theological reasoning.
This presentation follows John Marenbon's two complete chapters in the
Fordham manuscript of \emph{Mediaeval and Renaissance Logic}:
\emph{The Latin Tradition of Logic to 1100} and
\emph{Logic at the Turn of the Twelfth Century}
\cite[pp.~1--63]{Marenbon2007FordhamEarlyLatin}
\cite[pp.~65--81]{Marenbon2007FordhamTurnTwelfth}.
These manuscript versions are distinct from the final 2008 published chapters.

Marenbon's first chapter presents an attributed historical argument.
His second explicitly offers a guide to surviving material, dating,
interpretive difficulties and research questions, not a settled doctrine
of every early-twelfth-century school. The local formal comparisons below
make their assumptions explicit. They do not turn a historical term logic
into modern propositional calculus, nor credit cited medieval editions
as independently established originals. The separate Wilks chapter on
Abelard and his contemporaries is not replaced by these chapters'
references to Abelard.

\section{A Latin corpus, not just three Greek books}
\label{sec:p13-early-latin-corpus}

In the narrow sense, \emph{logica vetus} names Porphyry's \emph{Isagoge}
and Aristotle's \emph{Categories} and \emph{On Interpretation} in translation.
By the eleventh century the working curriculum also included Boethius's
categorical syllogisms, hypothetical syllogisms and topical argument.
His commentaries supplied ways to read the translated texts.
Marius Victorinus's treatise on definition, often attributed to Boethius,
Cicero's \emph{Topics}, and Boethius's commentary on it also mattered.
The \emph{Ten Categories}, Apuleius and the encyclopaedists had been
particularly important before Boethian teaching became dominant.

Marenbon's historical thesis concerns this changing corpus and its
institutional position. Logic became central to education beyond
elementary Latin grammar. Its importance cannot be measured only
by counting newly invented inference forms. In this period
\emph{dialectica} often names the whole subject, though a distinction
between demonstrative logic and persuasive dialectic could be drawn.
Marenbon acknowledges a different scholarly interpretation of that
vocabulary rather than using an etymology to settle the doctrines
\cite[pp.~1--3]{Marenbon2007FordhamEarlyLatin}.

\subsection{Cicero, Apuleius, and the earlier Latin resources}

Cicero's \emph{Topics}, placed in 44 BCE, is oriented toward finding
arguments for an orator, especially in legal disputes. Its short list
of relations includes genus/species, likeness, cause/effect and comparison.
The heir-to-a-dilapidated-house case uses a comparison with an heir
who need not replace a deceased slave. That is a proposed persuasive
analogy, not a formal entailment from resemblance alone.
Extrinsic appeal to authority differs from these intrinsic relations.

Cicero's seven modes derive, in an altered transmission, from Stoic
indemonstrables. Marenbon warns that definitions and examples do not
uniformly match. In particular the seventh mode, rendered with modern
propositional letters, is
\[
 \sim(p\mathbin{\&}q),\quad\sim p;\quad\text{therefore }q.
\]
The false assignment to both atoms is a counterexample.
This error is different from a later changed third-mode expression
that superficially looks related.

\begin{hblemma}{thm:p13-early-latin-modes}
 {proved=annex:p13.early-latin.modes}{Three negation patterns and their actual force}
In classical material-propositional semantics,
\(\sim(p\mathbin{\&}q),\sim p\) does not entail \(q\);
\(\sim(p\mathbin{\&}\sim q),p\) does entail \(q\);
and \(\sim(p\to\sim q),p\) also entails \(q\).
The last premise is equivalent to \(p\mathbin{\&}q\), not to a general
rule that \(p\) excludes \(\sim q\) across all situations.
\end{hblemma}

Apuleius's \emph{Peri hermeneias} treats non-modal Aristotelian
syllogistic, opposition and reduction to the first figure. Marenbon
notes its inference formulation, ordering and Latin terminology.
Its criticism of the day/light conditional confuses the embedded
content with the conditional operator: the occurrence of light in
a conditional need not change the meaning of light in the conclusion.
This criticism does not itself give a countermodel to correctly typed
conditional detachment.

Victorinus places logic within rhetoric, bringing the introductory Greek
texts toward Ciceronian Topics. His surviving definition work distinguishes
philosophical definition through predicables from rhetorical and other
definitions, including descriptions through parts. His lost works and
contested translations are not reconstructed from their mere mention.
The five predicables used in the philosophical account are genus,
species, differentia, proprium and accident. A definition through parts
can provide understanding without giving the same account of what
the thing is as a genus-and-differentia definition.
The \emph{Ten Categories}, once taken as Augustinian, belongs to an
unknown author associated with Themistius's circle. It simplifies
Aristotle for beginners and treats discourse about things and words
together rather than resolving every metaphysical question.

Martianus Capella's logical Book IV combines introductory category material,
Apuleian syllogistic and a seven-mode conditional list; Book V places
Topics within rhetoric. He sometimes retains Stoic numerical placeholders
for propositions. A source chain through Victorinus is a historical
interpretation, not an independently proved stemma
\cite[pp.~3--6]{Marenbon2007FordhamEarlyLatin}.

\section{Boethius: chosen interpretations and kinds of argument}
\label{sec:p13-early-latin-boethius}

\subsection{Translation and commentary are different contributions}

Marenbon presents Boethius as a selective transmitter, not either an
inventor of every doctrine or a mechanical translator of marginalia.
He translates the Organon except the \emph{Posterior Analytics},
but the earlier medieval curriculum initially uses only part of that
translation programme. Literal translations preserve interpretive
problems; the usual composite \emph{Categories} translation has its
own textual history. Shiel's marginalia explanation and the objections
to it remain scholarly alternatives, not a claim that every surviving
commentary has one known Greek source.

Porphyry is Boethius's preferred guide. The different-objects interpretation
reconciles Plato's intelligible objects with Aristotle's treatment of
the sensible world. Other Neoplatonic commentators use more explicitly
Platonic approaches to the logical texts. Porphyry's account of first
and second imposition distinguishes words for things from words for words.
Boethius's commentary on the \emph{Categories} develops that distinction
and a story of naming; the promised further commentary is not securely
established as an extant work
\cite[pp.~7--10]{Marenbon2007FordhamEarlyLatin}.

In \emph{On Interpretation}, utterance, writing, mental passions and
things form connected but distinct levels. Marenbon reads Boethius
as using a mental-discourse account: simple thoughts combine into
truth-evaluable complexes. The sameness of two mental concepts is
explained through their originating objects, not merely through the
physical resemblance of spoken word tokens.
This causal account of concepts is not identical to conventional
imposition of spoken names.

Boethius's interpretation of future contingents retains bivalence
while distinguishing definite from indefinite truth. Marenbon's
reading of the unqualified future assertion and its contingent
qualification is explicitly contrasted with other scholarship.
It is not a three-valued matrix or a probability assignment.
The later \emph{Consolation}'s prescience treatment is another solution,
not automatically the same account as the logical commentary
\cite[pp.~11--12]{Marenbon2007FordhamEarlyLatin}.

\subsection{The universals argument and the premise abstraction rejects}

Porphyry asks whether genera and species exist or occur only in thought;
if they exist, whether they are corporeal, and whether they are separate
or in sensible things. He declines to answer in the introductory text.
Marenbon reconstructs Boethius's second-commentary objection and response.
For one candidate universal, use the following explicitly modern
propositional abbreviations: \(E\), real existence; \(N\), numerical
oneness; \(M\), simultaneous commonness to many; \(T\), formation in
thought alone; \(C\), correspondence to the object just as it is;
and \(F\), being empty and false.

The source's argument has eleven stages:

\begin{enumerate}
\item Existence or formation in thought alone: \(E\lor T\).
\item What exists is numerically one: \(E\to N\).
\item What is common simultaneously to many is not numerically one:
  \(M\to\sim N\).
\item Genera and species are common in that way: \(M\).
\item Hence they do not exist in the proposed sense: \(\sim E\).
\item The remaining disjunct is formation in thought: \(T\).
\item A strictly corresponding universal thought would imply the proposed
  real existence: \(C\to E\).
\item The proposed real existence has already been rejected.
\item Hence these are not strictly corresponding thoughts: \(\sim C\).
\item The opponent adds that every non-corresponding thought is empty
  and false: \(\sim C\to F\).
\item Inquiry is consequently to be abandoned.
\end{enumerate}

\begin{hblemma}{thm:p13-early-latin-universals}
 {proved=annex:p13.early-latin.universals}{What the objection establishes without its final bridge}
The first four abbreviated premises and \(C\to E\) entail
\(\sim E,T,\sim C\), but not \(F\).
Adding \(\sim C\to F\) gives \(F\). A prescription to cease inquiry
requires a further normative bridge; it is not a truth-functional
consequence of the first nine stages alone.
\end{hblemma}

Boethius, on Marenbon's account, accepts the path through the ninth
stage and rejects the tenth. A line considered apart from a body
need not thereby be falsely thought. Abstraction can omit the mode
of concrete existence without being an empty error.
Marenbon distinguishes two solutions in the commentary.
In neutral abstractionism, resembling natures in particulars are
collected by thought; in realist abstractionism, singularity and
universality are grasped in one subject under different ways of access.
He observes that Boethius does not clearly mark the difference.
These interpretations must not be collapsed into one settled modern
ontology \cite[pp.~12--14]{Marenbon2007FordhamEarlyLatin}.

\subsection{Division, categorical syllogism, and qualified conversion}

\emph{On Division} distinguishes accidental division from division
of a whole into parts, a genus into species, and a word into meanings.
Whole/part treatment later provides material not equally available
in the other introductory books. The categorical monographs differ
in scope: \emph{On the Categorical Syllogism} has preliminaries and
syllogistic proper; the \emph{Introduction to Categorical Syllogisms}
develops the preliminaries in greater detail and has a disputed
unfinished or transmitted state.

The orientation of predication matters. A being predicated of every
B and B of every C yields A of every C. Changing the language to
the copula is, while retaining an older premise order, can obscure
the visual transparency without changing the valid inclusion.
Boethius's letter schemata and his ordinary examples are not treated
as identical phrasing.

\begin{hblemma}{thm:p13-early-latin-categorical}
 {proved=annex:p13.early-latin.categorical}{Predication chain and the conversion condition}
In the declared extension model, \(C\subseteq B\subseteq A\)
entails \(C\subseteq A\). \(A\subseteq B\) does not generally entail
\(B\subseteq A\). A further coextension condition, as in the
stated proprium example, licenses the reverse inclusion.
\end{hblemma}

The capacity-for-laughter example is a historically assumed proprium
of humans, not a newly established biological law. Boethius's treatment
of infinite names such as not-man and not-animal requires its own
term-negation interpretation; the complete original conversion tables
are not supplied by this chapter's report of their existence
\cite[pp.~14--15]{Marenbon2007FordhamEarlyLatin}.

\subsection{Hypothetical syllogisms are not silently propositional logic}

Boethius considers compounds using exclusive \emph{aut} and conditional
\emph{si}/\emph{cum}. He distinguishes accidental conditionals from
natural consequences. The human/animal example supplies relevant
inclusion without the causal direction of the earth-interposition/
lunar-eclipse example. The hot-fire/round-sky example concerns truths
without the same natural connection.
These are interpretations of his examples, not an empirical claim
that all modern causal instances fit a two-valued conditional.

The tabulation includes major premises with two, three or four term
positions; quality changes and modal or quantitative additions multiply
the cases. Marenbon reports Boethius's decision to develop the
non-modal, non-quantified combinations. He does not reproduce
every original Book II/III tabulation.
The letters in the hypothetical syllogisms denote predicate terms
in forms such as it is A, not simply arbitrary propositions.

For the major premise if not A then B, the ordinary accepted
minor cases give B from not A, and A from not B.
Inferring not A from B, or not B from A, requires more.
The additional assumption is that A and B are immediate contraries,
exhaustive and exclusive within their range of applicability.

\begin{hblemma}{thm:p13-early-latin-contraries}
 {proved=annex:p13.early-latin.contraries}{The precise extra force of immediate contrariety}
From \(\forall x(\sim Ax\to Bx)\), \(\sim Aa\) entails \(Ba\)
and \(\sim Ba\) entails \(Aa\) in classical semantics.
Neither \(Ba\) entails \(\sim Aa\) nor \(Aa\) entails \(\sim Ba\)
from that universal alone. If \(A,B\) partition the relevant domain,
both additional inferences become valid.
\end{hblemma}

Boethius's altered third mode illustrates the translation risk.
The historical reading takes A and not-B to be incompatible.
The literal material-propositional expression
\(\sim(p\to\sim q),p\) is valid, as the earlier lemma proves;
its validity does not show that it expresses that historical
incompatibility. In a term model,
\(\sim\forall x(Ax\to\sim Bx)\) merely supplies an A-and-B
witness somewhere, not necessarily at the subject \(a\).
By contrast, universal incompatibility of A with not-B, together
with \(Aa\), gives \(Ba\).
Negation of a whole conditional, denial of a term connection,
and affirmation of a complementary predicate must remain distinct.

Marenbon follows Martin in criticizing a reading that treats Boethius
as simply devising a strange propositional negation. The literal
state-of-affairs account of destroying a conditional can be
compatible with a counterexample to its term relation, even when
Boethius's own explanatory resources are term-based.
The present classical comparison is not an attribution of classical
propositional syntax to Boethius
\cite[pp.~15--18]{Marenbon2007FordhamEarlyLatin}.

\subsection{Topical discovery and its maximal warrant}

\emph{De topicis differentiis} and the commentary on Cicero compare
and organize topical differentiae. Their role is discovery:
a question concerning whether A is B suggests a search for a
middle term related to both. Part/whole, genus/species, sign/signified
and comparison guide the search in an arguer's existing knowledge.
This is not a supplied complete procedure that produces a sound
argument from every desired conclusion.

In the providence example, the world supplies the whole and human
affairs the part. The two explicit premises are that the world is
governed by providence and that human affairs belong to it.
The conclusion additionally relies on a maximal proposition
transferring the relevant fitting property from whole to part.
The differentia suggests the middle; the maximal proposition
supplies the warrant. Adding the warrant as a premise can make
the argument formally valid, but does not establish the warrant's truth.

\begin{hblemma}{thm:p13-early-latin-topical-warrant}
 {proved=annex:p13.early-latin.topical-warrant}{All premises of the whole-to-part argument}
For declared relation \(\operatorname{Part}\) and property \(P\),
\(P(w)\), \(\operatorname{Part}(h,w)\), and
\(\forall x\forall y((\operatorname{Part}(x,y)\land P(y))\to P(x))\)
entail \(P(h)\). The two ground premises alone do not.
\end{hblemma}

Some maximal propositions express definitions or familiar deductive
patterns; others are broad defeasible generalizations.
The authority differentia does not convert agreement of all,
many or learned speakers into an analytic truth.
Marenbon's assessment of the quality of these warrants is part
of his philosophical historical reading, not a proof that every
topical argument has equal force
\cite[pp.~18--20]{Marenbon2007FordhamEarlyLatin}.

\section{Encyclopaedic, Carolingian, and ninth-century approaches}
\label{sec:p13-early-latin-carolingian}

The encyclopaedic and Boethian traditions overlap without being identical.
Cassiodorus's \emph{Institutiones} combine introductory Greek material,
Apuleian syllogistic, Victorinian definitions and Topics.
Its alternative author-attribution lists assign materially different
roles to Victorinus and Boethius. These manuscript differences
are not resolved merely by citing the usual printed list.
Isidore largely compiles this material, at times introducing errors.
The emphasis on Categories, and Martianus's placement of Topics in
rhetoric, also prevents a single account of an exclusively rhetorical
encyclopaedic logic \cite[pp.~21--22]{Marenbon2007FordhamEarlyLatin}.

\subsection{Alcuin and Theodulf: educational and argumentative uses}

Alcuin's \emph{Dialectica}, conventionally dated 786--790 in
Marenbon's reconsidered chronology, arranges borrowed material for
a Christian educational programme. Its unusually large use of the
\emph{Ten Categories} directs attention to understanding the sensible
world and the distinction between creation and God.
This differs from a curriculum centered on deriving syllogisms.
Its setting as a dialogue with Charlemagne gives that curriculum
a political as well as pedagogical function.

Theodulf's argument in the \emph{Libri Carolini}, written in the
official response to Nicaea, disputes the supposed identity of
kissing and adoring. His display uses Boethian interpretation and
Apuleian opposition, but Marenbon observes that elaborate syllogizing
is not what establishes that the two expressions differ.
The educational, theological and controversial uses of logic converge
without becoming the same argument method.
Marenbon revises his own earlier priority attribution and explains
why the relative chronology of Alcuin and Theodulf cannot be
settled from intellectual affinity alone.

The manuscript associated with Leidrad and passages associated with
Candidus Wizo show different combinations of Categories, commentary
and argument resources. Their availability does not prove that
every master used every work
\cite[pp.~22--25]{Marenbon2007FordhamEarlyLatin}.

\subsection{Macarius, Ratramnus, and different orders of dependence}

The dispute reported by Ratramnus arises from the question whether
human souls are one, many, or one-and-many.
Macarius and his pupil treat the species as prior, with particulars
depending on it and differentiated through accidents.
Ratramnus makes particulars the subjects that sustain accidents
and construes universals apart from particulars as thought.
Their competing readings of Boethius's subsisting and being-a-subject
language are not merely competing pronunciations of one thesis.
Marenbon qualifies the older labels Platonic or hyper-realist in
favour of an immanent-realism interpretation. The incomplete
surviving report remains a limit on reconstruction
\cite[pp.~26--27]{Marenbon2007FordhamEarlyLatin}.

\subsection{Eriugena's Categories and the danger of treating a metaphysics as a calculus}

The Martianus glosses associated with Eriugena distinguish universal
\emph{ousia} from the individual hypostasis receiving accidents.
They treat quality through participation rather than solely as an
Aristotelian accidental form. The gloss tradition has multiple
redactions; parallel ideas do not establish that every gloss is
an exact direct product of one author.

In the \emph{Periphyseon}, nature is divided into uncreated/creating,
created/creating, created/not-creating and uncreated/not-creating.
Its extended Category analysis uses Greek Christian Neoplatonic
sources as well as Latin logical material.
Marenbon carefully distinguishes evidence of actual use of the
\emph{Ten Categories}, Martianus and topical material from weaker
evidence concerning the \emph{Isagoge} and other Boethian works.

Eriugena's Category groupings are not a partition used in one uniform
formal semantics. The within/around-\emph{ousia} distinction places
quality, relation, condition, acting and being-acted-on within it,
yet allows certain roles to cross the proposed divisions.
Quantity, place, time and position supply the surrounding limits
in the first scheme; this does not prevent quality, condition or
relation from also occurring in the description of other Categories.
The rest/motion grouping puts ousia, quantity, position and place
at rest and quality, relation, condition, time, acting and passion
in motion. Further reflections place universal essence above these
groupings and consider possible and impossible things.
This is not a supplied modern possible-world calculus.

The tree of genera and species is interpreted as creative procession,
not only as a classification of predicates. Universal essence and
particular substance are related under that assumption.
Bodies are distinguished from ousia: they arise through specified
incorporeal accidents together with essence; geometrical designation
and a natural body have different conditions.
Marenbon presents a charitable coherent line within those commitments
while retaining unresolved tensions. No set-theoretic inclusion
theorem establishes the metaphysical or theological claims
\cite[pp.~27--34]{Marenbon2007FordhamEarlyLatin}.

\section{Glosses and the new Boethian teaching}
\label{sec:p13-early-latin-glosses}

A gloss tradition is a family of overlapping annotations with additions,
omissions and different redactions, not one book by an author identified
from every shared sentence. Martianus, the \emph{Ten Categories} and
the \emph{Isagoge} generate different traditions.
Remigius's treatment can transmit and invert Eriugena's ousia/
hypostasis distinctions. The standard \emph{Ten Categories} glosses
increasingly explain the text's elementary logical doctrine rather
than pursue the Eriugenian theological project.
The special St Gall 274 material also uses Boethius's commentary.
The name once read as Icpa in an Isagoge explicit was reinterpreted
as Israel; this does not turn its composite quotations into an
independently invented semantic system.

Isolated \emph{On Interpretation} glosses include an unusually clear
temporal account of modality: necessary as always, impossible as
never, and contingent as varying between being and not-being.
That temporal-statistical interpretation is not automatically
necessity across every alternative course of events.
Distinguishing epistemic uncertainty about a past fact from openness
of a future event will matter again below
\cite[pp.~34--37]{Marenbon2007FordhamEarlyLatin}.

Around the late tenth century, fuller use of Boethius's monographs
and commentaries changes the curriculum. Marenbon distinguishes
surviving-copy evidence from the hypothesis of how the source
collection reached the schools. The decline in use of the
\emph{Ten Categories} and growth of the Boethian Categories material
are evidence for changed teaching, not a proof of every lost lecture.

Notker combines vernacular translation of Aristotle with older
Apuleian and Ciceronian resources in the syllogistic work attributed
to him. His elaboration of the altered third mode interprets
day and light through a natural connection and the exclusion of
opposites. Gerbert's curriculum, as reported by Richer, includes
the Boethian logical books, while his \emph{De rationale et ratione uti}
addresses the ambiguity between being rational and actually exercising
reason. A disposition need not be a continuously occurring act.
Abbo compares sources instead of merely copying them; his account
reduces hypothetical forms to seven general modes and distinguishes
Latin and Greek attitudes to immediate contraries.
The reported reduction and the disputed source of his list are
not a reconstruction here of every original hypothetical table
\cite[pp.~38--42]{Marenbon2007FordhamEarlyLatin}.

\section{Eleventh-century arguments in theology and semantics}
\label{sec:p13-early-latin-eleventh}

\subsection{Berengar's continuity dilemma}

Marenbon treats Lanfranc and Berengar as different users of logical
resources, not simply as one anti-dialectician opposed to one
dialectician. Berengar distinguishes signification from substantial
change and imposes a continuity requirement on the disputed
sacramental account.

His dilemma has two branches. If replacement flesh has only accidents
resembling the bread's, numerical continuity of the original bread
has not been established. If the very same accidents persist after
the bread's subject is removed, that conflicts with his account
of accidents individuated through their subject.
Under his continuity and dependence assumptions, neither branch
delivers the opponent's required transition.
The first branch is regarded as metaphysically possible but insufficient;
the second as capable of satisfying continuity but impossible.
Similarity, numerical identity, and persistence of a dependent
accident are distinct premises of this attributed argument.
It is not a religious ruling or an independently proved general
metaphysics of matter \cite[pp.~43--45]{Marenbon2007FordhamEarlyLatin}.

\subsection{Damian, Anselm, and the scope of necessity}

Peter Damian distinguishes restoring a present physical state from
making it the case that a past event never occurred.
His discussion separates the necessity of a discourse condition
from the necessity of the event itself. In modern notation, these
are not equivalent:
\[
 \Box(p\to p),\qquad p\to\Box p.
\]
The first does not establish the second.
Marenbon then questions Damian's parallel treatment of past and future
and reconstructs the further appeal to an eternal present.
That appeal is a theological and temporal interpretation, not
a formal proof that humans can alter already occurred facts.

\begin{hblemma}{thm:p13-early-latin-modal-scope}
 {proved=annex:p13.early-latin.modal-scope}{Necessity of a consequence versus necessity of its event}
On a two-world reflexive Kripke frame, \(\Box(p\to p)\) can hold
while \(p\to\Box p\) fails.
Likewise \(\Box(F\to P)\) together with actual \(F\) entails actual
\(P\) at a reflexive world, but need not entail \(\Box P\).
Adding \(\Box F\) supplies the stronger conclusion.
\end{hblemma}

Anselm distinguishes preceding coercive necessity from sequent
necessity in the foreknowledge/prophecy problem.
Marenbon reads his alternating if, because and while formulations
as not clearly separating all the relevant modal scopes.
The temporal premise that earlier knowledge is now fixed is an
additional difficulty, not resolved merely by distinguishing
\(\Box(p\to p)\) from \(p\to\Box p\).
The theological example involving an earlier purification through
a later death is treated as the source's argumentative setting,
not a current doctrinal pronouncement.

Anselm's Category analysis in the \emph{Monologion} also uses different
senses of substance, individuality, place and time. A negative
result under one sense and a positive result under another are
not automatically a contradiction in one fixed predicate.
The example of gaining an age or similarity relation to a newly
born person without an intrinsic change illustrates the extrinsic
character of some relations.
Marenbon's reading allows this discussion to question how secure
the ordinary Category scheme is.

The \emph{Proslogion} and Gaunilo exchange have broader modern
interpretive disputes not resolved by this chapter.
The understanding of its central description draws on the
mental-word tradition. Holopainen's proposed connection of the
single argument with a topical middle is reported as an interpretation,
including Marenbon's doubt about taking a bare term alone as an
argument carrying conviction. The philosophical fragments have
their own scope, not a fully supplied modal system here
\cite[pp.~45--51]{Marenbon2007FordhamEarlyLatin}.

\subsection{De grammatico: denominatives and distinct semantic levels}

A denominative names through an accident: grammaticus through
knowledge of grammar, albus through whiteness. The thing named
and the substance bearing the quality are not automatically
the same category of semantic object.
The opening syllogism is:
every grammaticus is a man; every man is a substance;
therefore every grammaticus is a substance.
Its validity is not removed merely by observing that a
non-human rational grammar-knower could be possible.

The later solution distinguishes the word signifying a quality
from its appellating a substance. The albus case makes the context
dependence concrete: a white thing in a house need not identify
a man or horse, while pointing at horses supplies a different
contextual identification.
Per-se signification concerns what thought is caused independently
of that context; per-aliud signification uses contextual information.
Anselm's appellation appeals to ordinary use. Marenbon warns
against assuming it already has every property of a modern
reference relation.

\begin{hblemma}{thm:p13-early-latin-grammaticus}
 {proved=annex:p13.early-latin.grammaticus}{The syllogism, possibility boundary, and typed semantic coexistence}
In a classical first-order model,
\(\forall x(Gx\to Mx),\forall x(Mx\to Sx)\) entails
\(\forall x(Gx\to Sx)\).
A possible-world counterinstance to \(G\subseteq M\) need not
falsify that universal at the actual world.
A word's signifying a quality and appellating a substance can
coexist without one object being both quality and substance.
\end{hblemma}

Marenbon distinguishes two moves in the dialogue.
The narrowly Aristotelian reading takes the claim grammaticus
is a quality as a metalinguistic assertion that the word signifies
a quality, using the discussion of substance words as an analogy.
The earlier broader reading permits denomination to be considered
with respect to the person or with respect to knowledge of grammar.
These are different interpretations, not one new theorem that
every entity belongs to two Categories.

Henry's Le\'{s}niewskian formal reconstruction, Adams's teaching-focused
interpretation and Marenbon's questioning of Category orthodoxy
are not interchangeable accounts. The difficulty and exploratory
character of the dialogue remain visible.
Anselm's brief immanent-realist remarks and hostile description
of Roscelin do not independently supply Roscelin's complete
theory or all of Anselm's original works
\cite[pp.~51--57]{Marenbon2007FordhamEarlyLatin}.

\section{The turn of the twelfth century: reading a layered corpus}
\label{sec:p13-early-latin-turn-twelfth}

Marenbon's second chapter is deliberately not a survey assigning
one definitive system to each named master. Material traditionally
treated as preceding Abelard may overlap chronologically with
his writing, including possible 1120s or 1130s layers.
Around 1110, the curriculum listed by Abelard consists of the two
Aristotelian works, Porphyry and four Boethian books.
The Topics in that list means \emph{De topicis differentiis},
not Aristotle's own Topics. Evidence of earlier knowledge of
the \emph{Prior Analytics} is distinguished from its later serious
influence \cite[pp.~65--66]{Marenbon2007FordhamTurnTwelfth}.

\subsection{Commentary form and the limits of authorship attribution}

Literal commentaries emphasize close explanatory paraphrase.
Composite commentaries combine it with more discursive discussion.
Problem commentaries concentrate on difficulties.
These are overlapping formats, not completely disjoint classes.
Lecture records, prepared teaching texts, corrected notes and
later reuse can all contribute to the surviving layers.
In P3, for example, one version adds a chimaera-related reading
to the shared universal discussion. A single surviving manuscript
could conceal an analogous addition.

The source's principal catalogue groups are:

\begin{itemize}
\item Paris BN 13368: P5/H4/D7, traditionally linked with early Abelard,
  but with that attribution strongly questioned.
\item Munich clm 14779: P7/H5/B1; H5 may preserve a fuller lecture version
  associated with H4, without securely identifying the lecturer.
\item Oxford Corpus Christi 233: H7, related more distantly to H4/H5.
\item Pommersfelden 16/2764: P4a/P4b, B3/B26, SC6/SH6 and notes;
  the manuscript's early dating constrains this material.
\item Composite P3: three witnesses with an earliest-preserved version
  and two independent revisions.
\item C8/C7/C14: six witnesses and multiple revision lines;
  H11/H9, D8, B8/B10 and SH3 have their own version relations.
\end{itemize}

The identifiers are source catalogue labels, not proofs of common
authorship or complete equivalence. P6/C6 remain little studied
in the chapter's account. P14/P16 and the further composite
material are not automatically dated by their dependence on Boethius.
Shared prologues, repeated passages and several manuscripts can
indicate teaching transmission without yielding a unique author
\cite[pp.~66--70]{Marenbon2007FordhamTurnTwelfth}.

\subsection{Treatises, grammar, rhetoric, and the evidential triangle}

The Dialecticas of Abelard and Gerlandus reorganize teaching material
rather than abandoning commentary resources altogether.
Their dating and direction of influence are disputed.
Gerlandus is associated with Besançon, not simply identified
with an eleventh-century predecessor. The Introductiones
concentrate on constructing propositions and arguments,
especially topical argument; uncertain manuscript dates and
attributions limit claims about their original authors.
Odo's discussion of original sin contains logical material,
but its surviving scope is not substituted for every lost
logical treatise attributed to him.

The Priscian Glosulae, Notae Dunelmenses and rhetorical commentaries
bring semantic and universal debates across the trivium.
The redating of a manuscript changes a proposed terminus ante quem.
Master G's identification with William and his recorded differences
from the Glosulae give conditional dating evidence, not an
unqualified authorial signature.
Abelard reports his clash with William over universals in a
rhetorical teaching context; the topic does not belong exclusively
to one course called logic
\cite[pp.~70--73]{Marenbon2007FordhamTurnTwelfth}.

The two prominent testimonies need different qualifications.
Anselm's description of Roscelin is adversarial and may caricature
the logical position. Abelard's retrospective account of William's
material-essence and non-difference positions is written with
an autobiographical purpose. It does not establish that every
argument later used by Abelard was already given around 1108,
or that William invented the first position himself.
Chronicle references to in-voce and in-re teaching add evidence,
but names occurring in logical examples are only possible links
between a master and a teaching record
\cite[pp.~73--75]{Marenbon2007FordhamTurnTwelfth}.

\subsection{In voce is not automatically an ontology}

Reading genus, species or even individual-category expressions
as words can be an exegetical strategy. It does not alone entail
that only particulars exist, nor that a writer had a complete
nominalist semantics. Gerlandus's unresolved passages and
the uncertainty surrounding Roscelin remain real limits.
Iwakuma's proto-vocalist/vocalist distinctions and proposed
William corpus are presented together with Marenbon's
qualifications of his own earlier view.

At least three further reasoning problems remain central.
Utterances are physical events or objects and bearers of meaning,
so their placement under substance, quantity or quality cannot
be settled by ignoring either role.
Potential and actual signification raise questions about
significative expressions and empty terms such as chimaera.
Topical argument supplies both a theory of discovery and an
explanatory practice of identifying the warrant in an argument.
The chapter points to these tasks without giving a completed
semantic theory for every anonymous source.

This early-twelfth-century treatment therefore ends with a
qualified account of what the corpus makes available for reasoning,
not a fabricated author-by-author calculus. The complete
Abelard treatment in Wilks, full original medieval texts,
later versions and final published chapter comparisons retain
their separate scope
\cite[pp.~75--79]{Marenbon2007FordhamTurnTwelfth}.
